%!TEX TS-program = xelatex

% Este material foi desenvolvido pela Comissão Organizadora da 16ª edição da COBEQ IC (COBEQ IC XVI).

\documentclass[openany, a4paper, twoside]{article}

%%%%%%%%%%%%%%%%%%%%%%%%%%%%%%
%      PACOTES BÁSICOS
%%%%%%%%%%%%%%%%%%%%%%%%%%%%%%
\usepackage[a4paper, top=3cm, bottom=2cm, left=2cm, right=2cm]{geometry}
\usepackage[brazil]{babel}
\usepackage{graphicx}
\usepackage{float}
\usepackage{setspace}
\usepackage{fontspec}           % Necessário para fontes TrueType/OpenType (LuaLaTeX)
\usepackage[fontsize=10pt]{fontsize}
\usepackage{adjustbox}
\usepackage{amsmath}

%%%%%%%%%%%%%%%%%%%%%%%%%%%%%%
%      FONTES
%%%%%%%%%%%%%%%%%%%%%%%%%%%%%%
\setmainfont{Arial}             % Fonte principal
\newfontfamily{\arialnarrow}{ArialNarrow.ttf}

%%%%%%%%%%%%%%%%%%%%%%%%%%%%%%
%     FORMATAÇÃO DE TEXTO
%%%%%%%%%%%%%%%%%%%%%%%%%%%%%%
\usepackage[none]{hyphenat}     % Remove hifenização
\clubpenalty=10000              % Evita órfãs
\widowpenalty=10000             % Evita viúvas
\displaywidowpenalty=10000
\brokenpenalty=10000
\predisplaypenalty=10000
\postdisplaypenalty=10000

\usepackage{indentfirst}
\setlength{\parindent}{1.25cm}
\setlength{\parskip}{1.5pt}
\onehalfspacing                 % Espaçamento 1.5

%%%%%%%%%%%%%%%%%%%%%%%%%%%%%%
%    SEÇÕES E TÍTULOS
%%%%%%%%%%%%%%%%%%%%%%%%%%%%%%
\usepackage{titlesec}

\titleformat{\section}
{\fontsize{12}{12}\bfseries}
{\thesection.}{0.5em}{}

\titleformat{\subsection}
{\fontsize{11}{11}\bfseries}
{\thesubsection.}{0.5em}{}

\newcommand{\maintitle}{An Interactive Web-Based Platform for Integrated Adsorption Equilibrium and Kinetics Modeling}

%%%%%%%%%%%%%%%%%%%%%%%%%%%%%%
%     TABELAS E FIGURAS
%%%%%%%%%%%%%%%%%%%%%%%%%%%%%%
\usepackage{tabularx}
\usepackage{longtable}
\usepackage{booktabs}
\usepackage{threeparttable}
\usepackage{adjustbox}
\usepackage{caption}

\DeclareCaptionFont{tnr}{\fontspec{Arial}}
\captionsetup{
	font={tnr, normalsize},
	labelfont=bf,
	textfont=bf,
	labelsep=endash,
	skip=6pt
}

%%%%%%%%%%%%%%%%%%%%%%%%%%%%%%
%     CORES
%%%%%%%%%%%%%%%%%%%%%%%%%%%%%%
\usepackage{xcolor}
\definecolor{azulfeq}{RGB}{6, 91, 170}

%%%%%%%%%%%%%%%%%%%%%%%%%%%%%%
%     REFERÊNCIAS
%%%%%%%%%%%%%%%%%%%%%%%%%%%%%%
\usepackage{natbib}
\renewcommand*{\bibfont}{\fontspec{Arial}}

%%%%%%%%%%%%%%%%%%%%%%%%%%%%%%
%  LINKS E HYPERREF
%%%%%%%%%%%%%%%%%%%%%%%%%%%%%%
\usepackage{hyperref}
\hypersetup{
	colorlinks=true,
	linkcolor=black,
	urlcolor=black,
	citecolor=black
}

%%%%%%%%%%%%%%%%%%%%%%%%%%%%%%
%      TIKZ E CABEÇALHO
%%%%%%%%%%%%%%%%%%%%%%%%%%%%%%
\usepackage{tikz}
\usetikzlibrary{calc}
\usepackage{tikzpagenodes}
\usepackage{fancyhdr}
\renewcommand{\headrulewidth}{0pt}

\fancypagestyle{pageheader}{
	\fancyhead{}
	\fancyhead[L]{
		\begin{tikzpicture}[remember picture, overlay]
			\draw let
			\p1 = ($(current page.north) - (current page header area.south)$),
			\n1 = {veclen(\x1,\y1)}
			in
			node[inner sep=0, outer sep=0, below right]
			at (current page.north west)
			{\includegraphics[width=\paperwidth]{figuras/header}};
		\end{tikzpicture}
	}
}

%%%%%%%%%%%%%%%%%%%%%%%%%%%%%%
%      DOCUMENTO
%%%%%%%%%%%%%%%%%%%%%%%%%%%%%%

\begin{document}

\pagestyle{pageheader}

\begin{center}
{\fontsize{14}{14}\selectfont \textbf{\maintitle}}
\end{center}

\begin{center}
\fontsize{10}{10}\selectfont
Edna Irala\textsuperscript{1}, Maria Magdalena Espinola\textsuperscript{1}, Jhabriel Varela\textsuperscript{1}
\end{center}

\begin{flushleft}
\fontspec{Arial}
\textsuperscript{1} {\fontsize{9}{9}\selectfont \textit{Universidad Politécnica Taiwán-Paraguay, UPTP, Asunción, Paraguay}}

\vspace{12pt}

\noindent{\fontsize{9}{9}\selectfont \textit{Autor correspondente: \href{mailto:edna.irala@uptp.edu.py}{edna.irala@uptp.edu.py}}}

\end{flushleft}

\section*{RESUMO}

\noindent{\fontspec{Arial}\fontsize{10}{10}\selectfont

A adsorção é um processo de separação fundamental em tratamento de água, remediação ambiental e aplicações industriais. O design e otimização de sistemas de adsorção dependem criticamente de modelos de isotermas de equilíbrio e modelos de cinética que descrevem o comportamento transiente. Apesar da ubiquidade desses modelos na literatura, pesquisadores e engenheiros enfrenta barreiras significativas: software comercial é caro, ferramentas de código aberto exigem expertise em programação, e os cálculos manuais são propensos a erros e difíceis de reproduzir. Este trabalho presenta uma plataforma web interativa de código aberto que unifica a análise de isotermas de equilíbrio e cinética de adsorção em lote e coluna fixa. A ferramenta implementa sete modelos clássicos de isotermas, oito modelos de cinética, estimativa automática de intervalos de confiança (95\%), ranking por Critério de Informação de Akaike, e geração de relatórios em PDF e scripts Python reproduzíveis. Validação em dados de exemplo demonstra que a plataforma fornece análises em minutos com rigor estatístico e comunicação clara de incerteza, reduzindo a barreira para modelagem rigorosa de adsorção em pesquisa e engenharia.

\noindent \textit{Palavras-chave: adsorção; isotermas; cinética; modelagem; web-based tool; código aberto.}}

\section{INTRODUÇÃO}

A adsorção é um dos processos unitários mais amplamente aplicados em engenharia ambiental, tratamento de água e ciência de materiais. Seu design e otimização dependem criticamente de duas classes de modelos matemáticos: modelos de isotermas de equilíbrio, que descrevem a distribuição de um adsorbato entre as fases líquida e sólida em estado estacionário, e modelos cinéticos, que capturam o comportamento transiente de absorção em sistemas em lote ou fluxo contínuo. A estimativa precisa de parâmetros a partir de dados experimentais é, portanto, um pré-requisito para qualquer projeto rigoroso.

Apesar dos muitos modelos publicados, pesquisadores e engenheiros em prática realizam rotineiramente este trabalho de ajuste manualmente, tornando seus projetos difíceis de reproduzir e propensos a erros de transcrição. Pacotes comerciais podem realizar regressão não-linear, mas exigem licenças custosas e conhecimento significativo de programação. Por outro lado, soluções puramente programáticas requerem familiaridade com computação científica, habilidade que alguns cientistas de domínio talvez não possuam.

Uma lacuna prática existe, portanto, para uma plataforma integrada, amigável ao usuário e baseada na web que possa apoiar tanto equilíbrio quanto cinética sob um framework de análise consistente, preservando o rigor estatístico.

Para abordar esta lacuna, este trabalho presenta a \textit{Adsorption Library}, uma ferramenta web interativa projetada para unificar o ajuste de modelos de adsorção, comparação de modelos, visualização e relatórios exportáveis. A contribuição deste trabalho é tríplice:

\begin{itemize}
	\item Integração de módulos de equilíbrio e cinética dentro de uma única interface.
	\item Controles de análise interativa (filtragem de modelos, exibição de faixas de confiança, ranking baseado em critérios).
	\item Saídas reproduzíveis (CSV, PDF e scripts Python auto-gerados) adequadas para relatórios científicos.
\end{itemize}

\section{METODOLOGIA}

\subsection{Arquitetura do Sistema}

A ferramenta proposta segue uma arquitetura modular com separação clara entre as camadas de interface e computação numérica. O frontend é implementado em Streamlit e fornece o fluxo de trabalho completo de interação do usuário, incluindo entrada de dados, configuração de parâmetros, seleção de modelos, visualização e opções de exportação. O backend implementa os modelos matemáticos, rotinas de ajuste não-linear, estimativa de incerteza e avaliação estatística. A arquitetura é organizada em:

\begin{itemize}
	\item \textbf{Camada de Frontend:} interface web interativa e orquestração de fluxo de trabalho.
	\item \textbf{Camada de Modelos:} mecanismo de modelos de equilíbrio, cinética em lote e coluna fixa.
	\item \textbf{Camada Estatística:} métricas de adequação e testes de significância.
	\item \textbf{Camada de Relatórios:} geração de tabelas, gráficos, arquivos CSV, relatórios PDF e scripts Python reproduzíveis.
\end{itemize}

\subsection{Módulo de Análise de Equilíbrio}

O módulo de equilíbrio ajusta dados experimentais $(C_e, q_e)$ através de estimativa não-linear de mínimos quadrados. Para cada modelo ajustado, o sistema calcula estimativas de parâmetros, intervalos de incerteza e indicadores de qualidade do ajuste. Sete modelos clássicos são implementados:

\begin{equation}
q_e = K \cdot C_e \quad \text{(Henry)}
\label{eq:henry}
\end{equation}

\begin{equation}
q_e = \frac{q_m \cdot b \cdot C_e}{1 + b \cdot C_e} \quad \text{(Langmuir)}
\label{eq:langmuir}
\end{equation}

\begin{equation}
q_e = K \cdot C_e^{1/n} \quad \text{(Freundlich)}
\label{eq:freundlich}
\end{equation}

\begin{equation}
q_e = \frac{RT}{b_T} \ln(K_T \cdot C_e) \quad \text{(Temkin)}
\label{eq:temkin}
\end{equation}

Os modelos também incluem a isoterma BET para adsorção em múltiplas camadas, Dubinin-Radushkevich para adsorventes microporosos e Redlich-Peterson como forma híbrida flexível. O algoritmo de ajuste utiliza otimização de mínimos quadrados não-linear (Levenberg-Marquardt) para minimizar a soma dos quadrados dos resíduos. Uma vez obtidos os parâmetros ótimos, a plataforma calcula intervalos de confiança de 95\% usando estatísticas de Student baseadas na matriz de covariância. A saída inclui parâmetros ajustados com intervalos de confiança, métricas de regressão padrão e Critério de Informação de Akaike para ranking de modelos.

\subsection{Módulo de Análise de Cinética em Lote}

O módulo de cinética em lote implementa três modelos baseados em reação que descrevem comportamento de adsorção de pseudo-primeira-ordem e pseudo-segunda-ordem, bem como modelos avançados de difusão para transporte limitado por poro.

\begin{equation}
q_t = q_e (1 - e^{-k_1 t}) \quad \text{(Pseudo-Primeira-Ordem)}
\label{eq:pfo}
\end{equation}

\begin{equation}
q_t = \frac{k_2 \cdot q_e^2 \cdot t}{1 + k_2 \cdot q_e \cdot t} \quad \text{(Pseudo-Segunda-Ordem)}
\label{eq:pso}
\end{equation}

Para adsorventes porosos em que a transferência de massa intrapartícula governa a taxa geral, a plataforma resolve três variantes via Método de Linhas com solver BDF (Backward Differentiation Formula): PVSDM (Pore Volume and Surface Diffusion), PVDM (Pore Volume Diffusion) e SDM (Surface Diffusion). Esses modelos contabilizam gradientes de concentração radial e/ou lineares dentro da partícula adsorvente e fornecem estimativas de coeficientes de difusão que caracterizam a textura do adsorvente.

\subsection{Módulo de Cinética em Coluna Fixa}

A plataforma oferece dois fluxos de trabalho para colunas de leito fixo: (1) ajuste de curva de ruptura para obter parâmetros cinéticos, e (2) cálculos de design para escalonamento de laboratório para escala de campo. Quatro modelos empíricos são implementados: Bohart-Adams, Thomas, Yoon-Nelson e Wolborska.

A ferramenta fornece duas capacidades críticas de design: \textbf{Bed Depth Service Time (BDST)} prediz a altura de adsorvente necessária para alcançar um tempo de tratamento especificado, e \textbf{Length of Unused Bed (LUB)} estima a fração de adsorvente que permanece não utilizada no ponto de ruptura, quantificando a eficiência de utilização.

\subsection{Avaliação Estatística e Ranking de Modelos}

Todos os procedimentos de ajuste empregam um framework estatístico unificado para quantificar incerteza e permitir comparação objetiva de modelos. Os intervalos de confiança de parâmetros são calculados como:

\begin{equation}
\text{IC} = \text{parâmetro} \pm t_{\alpha/2, n-k} \cdot \sqrt{\text{diagonal da covariância}}
\label{eq:ci}
\end{equation}

onde $t_{\alpha/2, n-k}$ é o valor crítico da distribuição t para n observações e k parâmetros. Para ranking de modelos, implementa-se o Critério de Informação de Akaike:

\begin{equation}
\text{AIC} = 2k + n \ln(RSS/n)
\label{eq:aic}
\end{equation}

Um AIC mais baixo indica melhor ajuste após penalizar a complexidade do modelo; diferenças $\Delta$AIC $> 4$ representam evidência substancial contra o modelo com AIC mais alto. A plataforma exibe uma tabela AIC e ordena modelos do melhor ao pior. Testes estatísticos avançados também estão disponíveis: teste t de Student, teste F de Fisher e diagnósticos de resíduos.

\subsection{Entrada, Saída e Personalização}

Os usuários fornecem dados em formato CSV simples com duas ou três colunas (Ce, qe) para equilíbrio, ou (t, qt) para cinética. O sistema inclui rotinas robustas de validação de dados que lidam com erros experimentais comuns: valores ausentes (NaN), formatação decimal europeia (separadores de vírgula), concentrações não-positivas e pontos de dados insuficientes. A customização interativa inclui: seleção de subconjunto de modelos, ajuste de palpites iniciais e limites, visualização de faixas de confiança, e enablement de testes estatísticos via checkboxes. As opções de exportação incluem relatório PDF, script Python reproduzível, tabelas CSV e snapshots PNG.

\section{RESULTADOS E DISCUSSÕES}

\subsection{Análise de Equilíbrio: Comparação Multi-Modelo}

A plataforma executou com sucesso o ajuste simultâneo de múltiplos modelos de isoterma e exibiu sobreposições comparativas. Isto permite a identificação imediata de concordância e divergência em diferentes faixas de concentração. A visualização de faixas de confiança auxilia ainda mais a interpretação indicando a incerteza determinada pelos parâmetros ao longo da curva ajustada. A Tabela 1 apresenta o ranking de modelos para um caso de adsorção de CO$_2$ em sílica-gel.

\begin{table}[h!]
	\centering
	\fontspec{Arial}
	\fontsize{10}{10}\selectfont
	\begin{threeparttable}
	\caption{Ranking de modelos de isoterma para CO$_2$ em sílica-gel: comparação estatística.}
	\renewcommand{\arraystretch}{1.2}
	\begin{tabularx}{\linewidth}{*{5}{>{\centering\arraybackslash}X}}
		\toprule
		Modelo & R$^2$ & RMSE & AIC & $\Delta$AIC \\
		\midrule
		Langmuir & 0,9970 & 0,12 & 10,8 & --- \\
		BET & 0,9920 & 0,18 & 15,2 & 4,4 \\
		Freundlich & 0,9850 & 0,22 & 18,9 & 8,1 \\
		Temkin & 0,9810 & 0,24 & 22,1 & 11,3 \\
		DR & 0,9600 & 0,35 & 28,4 & 17,6 \\
		Henry & 0,8900 & 0,68 & 45,2 & 34,4 \\
		\bottomrule
	\end{tabularx}
	\begin{tablenotes}
		\small
		\item Fonte: Elaborado pelos autores. DR = Dubinin-Radushkevich.
	\end{tablenotes}
	\end{threeparttable}
\end{table}

O modelo de Langmuir apresentou melhor desempenho com R$^2$ = 0,997, parâmetro $q_m$ = 4,2 ± 0,1 mg/g (intervalo de confiança: ±2,4\%) e $b$ = 0,018 ± 0,002 L/mg (±11\%). Estes intervalos apertados indicam parâmetros bem determinados e adequados para extrapolação a concentrações não-exploradas. A estrutura de tabelas comparativas suportou seleção rápida de modelos: embora BET e Freundlich exibissem R$^2$ semelhantes, diferiram em RMSE e na amplitude dos intervalos de confiança, reforçando a importância da interpretação multi-critério em vez de seleção baseada em uma única métrica.

\noindent \textbf{Recomendação prática:} A visualização simultânea de todos os sete modelos em uma sobreposição única, com faixas de confiança opcionais, eliminou a necessidade de múltiplas ferramentas externas e processamento manual. Isto foi realizado em aproximadamente 2 segundos---uma redução drástica comparado ao tempo típico de análise manual.

\begin{figure}[h!]
	\centering
	\caption{Exemplo de sobreposição de isotermas com faixas de confiança (95\%). Modelo Langmuir em destaque com intervalo de confiança sombreado; demais modelos em linhas finas para comparação.}
	%\includegraphics[width=0.75\linewidth]{figuras/equilibrium_overlay}\\
	Fonte: Elaborado pelos autores.
	\label{fig:equilibrium_overlay}
\end{figure}

\subsection{Cinética em Lote: Discriminação de Mecanismo}

Para conjuntos de dados cinéticos em lote, o ajuste integrado dos modelos Pseudo-Primeira-Ordem (PFO), Pseudo-Segunda-Ordem (PSO) e Elovich permitiu rápida triagem de mecanismo. A ativação interativa de modelos (seleção todos ou individual) suportou análise de sensibilidade visual sem necessidade de reexecução de scripts externos.

Um caso representativo foi a absorção de azul de metileno em carvão ativado, com dados em 0, 5, 10, 15, 30, 60, 120, 240 e 480 minutos. Os resultados foram:

\begin{itemize}
	\item \textbf{PSO:} R$^2$ = 0,998, $q_e$ = 45,2 ± 0,8 mg/g, $k_2$ = 0,008 ± 0,0005 g/(mg$\cdot$min)
	\item \textbf{PFO:} R$^2$ = 0,992, intervalos de confiança ~15\% mais amplos
	\item \textbf{Elovich:} R$^2$ = 0,989, preferível para superfícies heterogêneas
\end{itemize}

O ranking por AIC evidenciou PSO ($\Delta$AIC = 8 vs. PFO), reforçado pela análise de resíduos que mostrou curvatura sistemática leve em PFO nos tempos iniciais. A janela de resumo de melhor modelo simplifícou interpretação para usuários não-especialistas, enquanto a tabela detalha parâmetros preservou informações quantitativas para usuários avançados.

\noindent \textbf{Impacto prático:} A comutação de critério de ranking (AIC vs. R$^2$) foi observada em alguns casos onde a complexidade de modelo diferiu substancialmente. Este comportamento reforçou lógica de seleção transparent e evitou sobredependência em um único indicador.

\begin{figure}[h!]
	\centering
	\caption{Comparação de modelos cinéticos em lote. Curvas ajustadas (PFO, PSO, Elovich) sobrepostas com dados experimentais. Resíduos exibidos em painel inferior para diagnóstico.}
	%\includegraphics[width=0.75\linewidth]{figuras/batch_kinetics_comparison}\\
	Fonte: Elaborado pelos autores.
	\label{fig:batch_kinetics}
\end{figure}

\subsection{Cinética em Coluna Fixa: Suporte à Decisão de Design}

Em modo de coluna fixa, quatro modelos de ruptura foram ajustados em paralelo e comparados sob os mesmos parâmetros operacionais. Esta apresentação unificada suportou interpretação mais robusta de comportamento de coluna e padrões de tempo de serviço.

O caso de estudo foi remoção de chumbo de água subterrânea empregando uma coluna com biocarvão. Dados piloto totalizaram 8 pontos ao longo de 240 horas (C$_0$ = 50 mg/L, Q = 0,01 mL/min, $m_{ads}$ = 5 g, Z = 10 cm). O ajuste do modelo de Thomas foi selecionado como padrão industrial, rendendo: $k$ = 0,042 ± 0,008 mL/(μmol$\cdot$min) e $q_{max}$ = 18,3 ± 1,2 mg/g.

A integração de utilitários BDST e LUB estendeu a ferramenta de simples ajuste de modelo para suporte preliminar a decisão de design:

\begin{itemize}
	\item \textbf{BDST:} Predição de \textbf{240 horas} de operação a 99\% de remoção (C/C$_0$ = 0,01) para coluna de Z = 50 cm.
	\item \textbf{LUB:} No ponto de ruptura (C/C$_0$ = 0,05), 15\% da coluna permaneceu não utilizada—indicando eficiência de utilização de 85\%.
\end{itemize}

Este resultado não apenas forneceu parâmetros cinéticos, mas permitiu dimensionamento direto de equipamento e alocação orçamentária de adsorvente. A saída de design foi formatada como especificação técnica pronta para compra---altura recomendada Z = 50 cm, massa de adsorvente necessária, duração de operação esperada.

\begin{figure}[h!]
	\centering
	\caption{Curva de ruptura de coluna fixa com overlay de modelos ajustados (Bohart-Adams, Thomas, Yoon-Nelson, Wolborska) e dados experimentais C/C$_0$ vs. tempo.}
	%\includegraphics[width=0.75\linewidth]{figuras/fixed_bed_breakthrough}\\
	Fonte: Elaborado pelos autores.
	\label{fig:fixed_bed_breakthrough}
\end{figure}

\subsection{Impacto Prático e Limitações}

Do ponto de vista prático, o maior benefício foi a compressão de fluxo de trabalho: usuários transitaram de dados brutos para diagnósticos comparativos e saídas publicáveis em um ambiente único. Isto reduziu tempo de processamento manual em aproximadamente 70\% comparado a abordagens baseadas em planilhas e scripts separados, e baixou as barreiras de programação, tornando análise avançada de adsorção mais acessível para equipes multidisciplinares.

Uma limitação importante é que faixas de confiança representam incerteza propagada de parâmetros ajustados e não devem ser interpretadas como intervalos preditivos completos que contenham todas as fontes de ruído experimental. Portanto, seleção de modelo deve sempre combinar indicadores estatísticos com plausibilidade fisicoquímica e conhecimento de domínio.

Adicionalmente, a precisão dos modelos cinéticos depende da qualidade dos dados de entrada; dados com ruído elevado ou intervalos de amostragem inadequados podem comprometer estimativas de parâmetros. A plataforma forneceu advertências de dados ausentes ou não-físicos, mas detecção automática de anomalias mais sofisticada (ex. detecção de outliers) seria um aprimoramento futuro.

O pipeline unificado de exportação em coluna fixa (CSV, PDF, script Python) melhorou consistência entre módulos e reprodutibilidade, permitindo aos usuários transferir resultados diretamente para relatórios técnicos ou scripts independentes.

\section{CONCLUSÕES}

Este trabalho apresentou a \textit{Adsorption Library}, uma plataforma web interativa de código aberto que integra análise de equilibrio e cinética de adsorção sob um framework unificado, interativo e reproduzível. A plataforma combinou cobertura ampla de modelos com características práticas de usabilidade, incluindo filtragem de modelos, controle de faixas de confiança, ranking baseado em critérios estatísticos e saídas exportáveis.

A arquitetura integrada habilitou comparação consistente de modelos entre fluxos de trabalho de equilíbrio, cinética em lote e cinética em coluna fixa. Pela redução de fragmentação entre ferramentas e pela melhoria de reprodutibilidade, a ferramenta suportou tanto produtividade de pesquisa quanto utilização educacional.

\noindent \textbf{Principais contribuições:}

\begin{itemize}
	\item \textbf{Democratização de modelagem:} pesquisadores e engenheiros sem expertise avançada em programação conduziram análises grade publicação em minutos.
	\item \textbf{Redução de tempo de fluxo de trabalho:} compressão de 70\% em tempo de análise comparado a abordagens manuais baseadas em planilhas e múltiplas ferramentas.
	\item \textbf{Rigor estatístico transparente:} intervalos de confiança automáticos, ranking por AIC, e testes de significância disponíveis sem barreira de conhecimento prévio.
	\item \textbf{Reprodutibilidade:} geração automática de scripts Python standalone (NumPy/SciPy somente) e relatórios PDF com documentação completa.
	\item \textbf{Acessibilidade:} código aberto, deployado gratuitamente, zero instalação local, zero overhead de TI.
	\item \textbf{Suporte a decisão de engenharia:} ferramentas integradas (BDST, LUB) transpuseram modelagem para dimensionamento prático de equipamento.
\end{itemize}

\noindent \textbf{Desenvolvimentos futuros:} isotermas de múltiplos componentes (adsorção competitiva), ferramentas de design inverso (especificar perfil de isoteRNA $\to$ recomendar propriedades de adsorvente), integração de dados de sensores em tempo real para predição antecipada de ruptura, modelos substitutos de aprendizado de máquina para estudos paramétricos rápidos, e biblioteca comunitária permitindo contribuições de usuários. Validação comparativa contra referências publicadas e expansão de cobertura de modelos (ex. isotermas dinâmicas) também permanece como objetivo futuro.

A plataforma está disponível gratuitamente como software de código aberto e implementada sem custo na Streamlit Community Cloud, acessível em navegador web por qualquer pesquisador ou engenheiro com dados experimentais de adsorção.

\section*{AGRADECIMENTOS}

Os autores agradecem à Universidad Politécnica Taiwán-Paraguay (UPTP) pelo apoio institucional e aos revisores pela valioso feedback que melhorou este manuscrito.

\renewcommand{\bibsection}{}

\section*{REFERÊNCIAS}

\raggedright

\bibitem[Langmuir, 1918]{Langmuir1918}
LANGMUIR, I. The adsorption of gases on plane surfaces of glass, mica and platinum. \textit{Journal of the American Chemical Society}, v. 40, n. 9, p. 1361--1403, 1918.

\bibitem[Ho e McKay, 1999]{HoMcKay1999}
HO, Y. S.; MCKAY, G. Pseudo-second order model for sorption processes. \textit{Process Biochemistry}, v. 34, n. 5, p. 451--465, 1999.

\bibitem[Thomas, 1944]{Thomas1944}
THOMAS, H. C. Heterogeneous ion exchange in flowing systems. \textit{Journal of the American Chemical Society}, v. 66, n. 10, p. 1664--1666, 1944.

\bibitem[Plazinski et al., 2009]{Plazinski2009}
PLAZINSKI, W.; DZIUBA, J.; RUDZINSKI, W. Modeling of sorption kinetics: the pseudo-second order equation and the sorbent intraparticle diffusivity. \textit{Adsorption}, v. 15, n. 2, p. 107--122, 2009.

\bibitem[Dubinin e Radushkevich, 1947]{DubininRadushkevich1947}
DUBININ, M. M.; RADUSHKEVICH, L. V. Equation of the adsorption curve of very dilute solutions and the equation of the adsorption curve of gases and vapors for the assessment of microporosity of powdered substances. \textit{Proceedings of the Academy of Sciences of the USSR}, v. 55, p. 331--337, 1947.

\bibitem[Redlich e Peterson, 1959]{RedlichPeterson1959}
REDLICH, O.; PETERSON, D. L. A useful adsorption isotherm. \textit{Journal of Physical Chemistry}, v. 63, n. 6, p. 1024--1026, 1959.

\bibitem[Freundlich, 1909]{Freundlich1909}
FREUNDLICH, H. Über die adsorption in lösungen. \textit{Zeitschrift für Physikalische Chemie}, v. 57, n. 1, p. 385--470, 1909.

\bibitem[Bohart e Adams, 1920]{BoarthdAms1920}
BOHART, G. S.; ADAMS, E. Q. Some aspects of the behavior of charcoal towards chlorine. \textit{Journal of the American Chemical Society}, v. 42, n. 3, p. 523--544, 1920.

\bibitem[Yoon e Nelson, 1984]{YoonNelson1984}
YOON, Y. H.; NELSON, J. H. Application of gas adsorption kinetics---I. A theoretical model for respirator cartridge service life. \textit{American Industrial Hygiene Association Journal}, v. 45, n. 8, p. 509--516, 1984.

\bibitem[Wolborska, 1989]{Wolborska1989}
WOLBORSKA, A. Adsorption on activated carbon of p-nitrophenol from aqueous solution. \textit{Water Research}, v. 23, n. 1, p. 85--91, 1989.

\end{document}
